\chapter{Análise descendente recursiva}

\section{Introdução}

Neste capítulo vamos estudar a primeira técnica para construção de
analisadores sintáticos: a análise descendente recursiva. Essa
técnica permite a construção de um analisador sintático como um
conjunto de funções, uma para cada não terminal da gramática da
linguagem considerada.

Em linguagens funcionais, como Haskell, uma abordagem elegante para
construção de analisadores descendentes recursivos utiliza os
chamados \textbf{combinadores de parsing}, que consistem de um
funções para construir analisadores simples e outras funções para
combinar analisadores existentes. Essa técnica permite descrever a
gramática da linguagem de forma direta. O código discutido neste
capítulo está disponível nos arquivos
\begin{flushleft}
  \verb|Parsing/Descendent/Combinators.hs|,\\
  \verb|Parsing/Descendent/Expr.hs| e\\
  \verb|Parsing/Descendent/ExpExample.hs|.
\end{flushleft}

\section{Combinadores de parsing}

A ideia central por trás dos combinadores de parsing é tratar um
analisador sintático como um \textbf{valor de primeira classe}: um
analisador é simplesmente uma função que recebe uma lista de símbolos
e retorna os possíveis resultados da análise. Essa representação
permite combinar analisadores simples para construir analisadores
mais complexos usando funções de ordem superior.

\subsection{O tipo Parser}

O tipo fundamental da biblioteca é definido como:

\begin{minted}{haskell}
newtype Parser s a =
  Parser { runParser :: [s] -> [(a,[s])] }
\end{minted}

Um valor do tipo \haskell{Parser s a} representa um
analisador que:
\begin{itemize}
  \item Consome símbolos do tipo \haskell{s} (por exemplo,
    \haskell{Char} ou algum tipo de token);
  \item Produz resultados do tipo \haskell{a} (por exemplo, uma
    árvore de sintaxe abstrata).
\end{itemize}

A função \haskell{runParser} extrai a função
subjacente, que recebe a entrada como uma lista de símbolos
\haskell{[s]} e devolve uma lista de pares
\haskell{[(a, [s])]}. Cada par contém um possível
resultado da análise e o restante da entrada ainda não consumida.

O uso de uma lista de resultados é o que torna esta representação
poderosa: a lista vazia, \haskell{[]}, sinaliza
\textbf{falha} (nenhum resultado possível), enquanto uma lista com
múltiplos elementos representa um analisador \textbf{ambíguo}, com
mais de uma interpretação possível para a entrada. Um analisador
determinístico retornará sempre no máximo um elemento na lista.

\subsection{Instâncias de classes de tipos}

A estrutura algébrica do tipo \haskell{Parser} é
rica: ele pode ser instanciado em várias classes de tipos de Haskell,
o que nos dá gratuitamente um conjunto de combinadores expressivos.

\textbf{Functor.} Um parser é um functor: dado um parser que produz
valores do tipo \haskell{a} e uma função
\haskell{f :: a -> b}, podemos construir um parser
que produz valores do tipo \haskell{b}, aplicando
\haskell{f} a cada resultado.

\begin{minted}{haskell}
instance Functor (Parser s) where
  fmap f p
    = Parser (\ s -> [(f x, s') |
                      (x,s') <- runParser p s])
\end{minted}

\textbf{Applicative.} Um parser é um functor aplicativo. O combinador
\haskell{pure} cria um functor que sempre tem sucesso
sem consumir entrada, retornando o valor fornecido. O operador
\haskell{<*>} sequencia dois parsers: o primeiro
produz uma função e o segundo produz um argumento; os resultados são combinados.

\begin{minted}{haskell}
instance Applicative (Parser s) where
  pure  a = Parser (\ts -> [(a, ts)])
  p1 <*> p2 = Parser (\ s -> [(f x, s2) |
           (f, s1) <- runParser p1 s,
           (x, s2) <- runParser p2 s1])
\end{minted}

\textbf{Alternative.} A instância de
\haskell{Alternative} nos fornece o combinador de
\textbf{escolha}. O parser \haskell{empty} falha
sempre. O operador \haskell{<|>} tenta o primeiro
parser e o segundo parser. Porém, dada a lazy evaluation de Haskell,
o segundo parser só será realmente executado se o primeiro falhar.
Isso se deve devido a concatenação dos resultados de ambos os parsers.

\begin{minted}{haskell}
instance Alternative (Parser t) where
  empty = Parser (\ _ -> [])
  p1 <|> p2 = Parser
     (\ s -> runParser p1 s ++
             runParser p2 s)
\end{minted}

\textbf{Monad.} A instância de \haskell{Monad}
permite que o segundo parser dependa do resultado do primeiro.

\begin{minted}{haskell}
instance Monad (Parser t) where
  return =  pure
  p >>= f  = Parser (\ts -> concat [ runParser (f a) cs' |
                                     (a,cs') <- runParser p ts ])
\end{minted}

Com essa instância, é possível escrever analisadores em notação
\haskell{do}, tornando o código mais legível para
parsers com múltiplos estágios sequenciais.

\textbf{MonadPlus.} Essa instância espelha a instância de
\haskell{Alternative}, fornecendo
\haskell{mzero} (falha) e
\haskell{mplus} (escolha) para uso no contexto monádico.

\subsection{Combinadores primitivos}

A biblioteca oferece dois parsers primitivos que servem como base
para todos os demais.

A função \haskell{item} consome exatamente um símbolo
da entrada, falhando se a entrada estiver vazia:

\begin{minted}{haskell}
item :: Parser t t
item = Parser (\ ts ->
                 case ts of
                   []     -> []
                   (c:cs) -> [(c,cs)])
\end{minted}

O combinador \haskell{sat} consome um símbolo somente
se ele satisfaz um predicado:

\begin{minted}{haskell}
sat :: (t -> Bool) -> Parser t t
sat p = do
          t <- item
          if p t then return t else mzero
\end{minted}

A partir de \haskell{sat}, derivamos diretamente os
combinadores para reconhecer símbolos e sequências específicas:

\begin{minted}{haskell}
symbol :: Eq s => s -> Parser s s
symbol c = sat (c ==)

token :: Eq s => [s] -> Parser s [s]
token = mapM symbol
\end{minted}

\haskell{symbol c} consome exatamente o símbolo
\haskell{c}; \haskell{token ts}
consome exatamente a sequência \haskell{ts},
aplicando \haskell{symbol} a cada elemento.

\subsection{Combinadores de alto nível}

Utilizando as funções primitivas, podemos construir um conjunto de
combinadores de alto nível.

\textbf{Repetição.} A classe \haskell{Alternative}
fornece \haskell{many} (zero ou mais repetições) e
\haskell{many1} (uma ou mais):

\begin{minted}{haskell}
many1    :: Parser s a -> Parser s [a]
many1 p  =  list <$> p <*> many p
\end{minted}

\textbf{Opção.} O combinador \haskell{option} executa
um parser opcionalmente, retornando um valor padrão em caso de falha:

\begin{minted}{haskell}
option :: Parser s a -> a -> Parser s a
option p v = p <|> succeed v
\end{minted}

\textbf{Pack.} O combinador \haskell{pack} sequencia
três parsers, descartando o resultado do primeiro e do terceiro e
retornando apenas o resultado do segundo. É útil para reconhecer
construções entre delimitadores:

\begin{minted}{haskell}
pack :: Parser s a -> Parser s b -> Parser s c -> Parser s b
pack p r q  =  pi32 <$> p <*> r <*> q
\end{minted}

Por exemplo, \haskell{pack (symbol '(') p (symbol ')')}
reconhece \haskell{p} entre parênteses.

\textbf{Separadores.} O combinador \haskell{listOf}
reconhece uma lista de itens separados por um delimitador:

\begin{minted}{haskell}
listOf      :: Parser s a -> Parser s b -> Parser s [a]
listOf p s  =  list <$> p <*> many (pi22 <$> s <*> p)
\end{minted}

\textbf{Repetição gulosa.} O combinador
\haskell{determ} trunca a lista de resultados ao
primeiro, tornando o parser determinístico. Sobre ele,
\haskell{greedy} e \haskell{greedy1}
implementam versões determinísticas de \haskell{many}
e \haskell{many1}:

\begin{minted}{haskell}
determ  ::  Parser s b -> Parser s b
determ p = Parser (\ ts ->
                     case (runParser p ts) of
                       []    -> []
                       (x:_) -> [x] )

greedy, greedy1  ::  Parser s b -> Parser s [b]
greedy   =  determ . many
greedy1  =  determ . many1
\end{minted}

\subsection{Aplicações para tipos concretos}

Com os combinadores acima, a biblioteca define analisadores para
construções léxicas comuns sobre \haskell{Parser Char}:

\begin{minted}{haskell}
digit :: Parser Char Int
digit = f <$> sat isDigit
  where
    f c = ord c - ord '0'

natural  :: Parser Char Int
natural  =  foldl (\a b -> a*10 + b) 0 <$> many1 digit

integer  ::  Parser Char Int
integer  =  (const negate <$> (symbol '-')) `option` id  <*>  natural

identifier :: Parser Char String
identifier =  list <$> sat isAlpha <*> greedy (sat isAlphaNum)

parens :: Parser Char a -> Parser Char a
parens p  =  pack (symbol '(') p (symbol ')')

spaces :: Parser Char String
spaces = greedy (sat isSpace)
\end{minted}

\subsection{Combinadores para encadeamento de operadores}

Um padrão muito comum em gramáticas de linguagens de programação é a
expressão com operadores infixos associativos. Para não-terminais da forma

\[
  E \to E \oplus T \mid T
\]

(associatividade à esquerda) ou

\[
  E \to T \oplus E \mid T
\]

(associatividade à direita), a biblioteca oferece os combinadores
\haskell{chainl} e \haskell{chainr}:

\begin{minted}{haskell}
chainr  ::  Parser s a -> Parser s (a -> a -> a) -> Parser s a
chainr pe po  =  h <$> many (j <$> pe <*> po) <*> pe
  where j x op  =  (x `op`)
        h fs x  =  foldr ($) x fs

chainl  ::  Parser s a -> Parser s (a -> a -> a) -> Parser s a
chainl pe po  =  h <$> pe <*> many (j <$> po <*> pe)
  where j op x  =  (`op` x)
        h x fs  =  foldl (flip ($)) x fs
\end{minted}

O argumento \haskell{pe} é o analisador para os
operandos (expressões de nível mais alto) e
\haskell{po} é o analisador para os operadores, que
devolve uma \textbf{função binária} a ser aplicada aos operandos.

Por exemplo, \haskell{chainl termParser opParser}
reconhece a gramática

\[
  E \to E + T \mid T
\]

produzindo uma árvore associativa à esquerda: a expressão
\haskell{1 + 2 + 3} é interpretada como
\haskell{(1 + 2) + 3}.

\section{O parser de expressões}

Para ilustrar o uso da biblioteca, vamos examinar o analisador
sintático para uma linguagem simples de expressões aritméticas com
variáveis. A linguagem é dada pela seguinte gramática:

\[
  \begin{array}{lcl}
    E & \to & E + T \mid T\\
    T & \to & T * F \mid F\\
    F & \to & n \mid x \mid (E)
  \end{array}
\]

onde \(n\) representa literais inteiros e \(x\) representa
identificadores. A implementação está dividida em dois estágios, como
é habitual: a \textbf{análise léxica}, que converte a string de
entrada em uma lista de tokens, e a \textbf{análise sintática}, que
consome os tokens e produz uma AST.

\subsection{Tokens e análise léxica}

Os tokens da linguagem de expressões são definidos pelo tipo:

\begin{minted}{haskell}
data Token
  = Id String
  | Number Int
  | Add
  | Mult
  | LParen
  | RParen
\end{minted}

O analisador léxico é construído com os próprios combinadores da
biblioteca. O lexer principal elimina espaços em branco e aplica
\haskell{tokenLexer} repetidamente:

\begin{minted}{haskell}
lexer' :: Parser Char [Token]
lexer' = (\_ x -> x) <$> spaces
                        <*> many ((\x _ -> x) <$> tokenLexer <*> spaces)
\end{minted}

O analisador \haskell{tokenLexer} tenta reconhecer,
em ordem, um identificador, um número ou um operador/delimitador.
Para os operadores e delimitadores, utiliza-se uma tabela de
associação entre caracteres e tokens:

\begin{minted}{haskell}
tokenLexer :: Parser Char Token
tokenLexer = (Id <$> identifier) <|> number <|> tokenlex

table :: [(Char, Token)]
table = [('+', Add), ('*', Mult),
         ('(', LParen), (')', RParen)]

tokenlex :: Parser Char Token
tokenlex = choice $ map (\(c,t) -> const t <$> symbol c) table
\end{minted}

O uso de \haskell{choice} com a tabela é um padrão
limpo: para cada par \haskell{(c, t)} na tabela,
construímos um parser que reconhece o caractere
\haskell{c} e retorna o token
\haskell{t}. \haskell{choice} combina
todos esses parsers com \haskell{<|>}.

A função \haskell{exprLexer} executa o lexer e extrai
o único resultado esperado:

\begin{minted}{haskell}
exprLexer :: String -> [Token]
exprLexer s
  = case runParser lexer' s of
      [] -> []
      ((x, "" ):_) -> x
      v -> error $ show v
\end{minted}

O padrão \haskell{(x, "")} garante que toda a entrada
foi consumida; caso haja entrada residual, um erro é sinalizado.

\subsection{A árvore de sintaxe abstrata}

A AST para as expressões é representada pelo tipo algébrico:

\begin{minted}{haskell}
data Expr
  = Var String
  | Lit Int
  | Expr :+: Expr
  | Expr :*: Expr
  deriving (Eq, Show)
\end{minted}

Os construtores \haskell{(:+:)} e
\haskell{(:*:)} são operadores infixos definidos no
tipo, o que torna o código do parser notavelmente próximo da gramática.

\subsection{O combinador gen para gramáticas com operadores}

Antes de apresentar o analisador sintático, é necessário introduzir o
combinador auxiliar \haskell{gen}, definido em
\verb|Parsing/Descendent/Expr.hs|:

\begin{minted}{haskell}
type Op s a = (s, a -> a -> a)

gen :: Eq s => [Op s a] -> Parser s a -> Parser s a
gen ops p = chainl p (choice (map f ops))
  where
    f (s,c) = const c <$> symbol s
\end{minted}

O tipo \haskell{Op s a} associa um símbolo
\haskell{s} (um token) a uma função binária que
combina dois valores do tipo \haskell{a}. O
combinador \haskell{gen} recebe uma lista de tais
associações e um parser para os operandos, e
constrói um parser para
expressões com os operadores dados, com
associatividade à esquerda.

Internamente, \haskell{gen} constrói o parser do
operador com \haskell{choice (map f ops)}: para cada
par \haskell{(s, c)}, cria um parser que consome o
token \haskell{s} e retorna a função
\haskell{c}. Em seguida, aplica
\haskell{chainl} para obter a associatividade à esquerda.

\subsection{O analisador sintático}

O analisador para expressões segue a gramática com três níveis de precedência:

\begin{minted}{haskell}
exprParser :: Parser Token Expr
exprParser = addtable `gen` termParser
  where
    addtable = [(Add, (:+:))]

termParser :: Parser Token Expr
termParser = multable `gen` factParser
  where
    multable = [(Mult, (:*:))]

factParser :: Parser Token Expr
factParser = numParser  <|>
             varParser  <|>
             parenExpr
  where
    parenExpr = pack lparen exprParser rparen
    lparen = symbol LParen
    rparen = symbol RParen
\end{minted}

A estrutura do parser espelha diretamente a gramática:

\begin{itemize}
  \item \haskell{exprParser} reconhece somas de termos
    ($E \to T + E \,|\, T$). Utiliza \haskell{gen} com a
    tabela de adição, tendo \haskell{termParser} como
    analisador dos operandos.
  \item \haskell{termParser} reconhece produtos de
    fatores ($T \to F * T \,|\, F$). Utiliza
    \haskell{gen} com a tabela de multiplicação,
    tendo \haskell{factParser} como analisador dos operandos.
  \item \haskell{factParser} reconhece fatores atômicos
    ($F \to n \,|\, x \,|\, (E)$): literais inteiros,
    variáveis ou expressões parentesadas.
\end{itemize}

Os analisadores para variáveis e números usam
\haskell{sat} para verificar o construtor do token:

\begin{minted}{haskell}
varParser :: Parser Token Expr
varParser = f <$> sat isVar
  where
    isVar (Id _) = True
    isVar _      = False
    f (Id v) = Var v
    f _ = error "impossible!"

numParser :: Parser Token Expr
numParser = f <$> sat isNum
  where
    isNum (Number _) = True
    isNum _ = False
    f (Number n) = Lit n
    f _          = error "impossible!"
\end{minted}

A função de entrada do parser converte a lista de tokens em um valor
\haskell{Either}, retornando a AST ou uma mensagem de erro:

\begin{minted}{haskell}
expr :: [Token] -> Either String Expr
expr ts
  = case runParser exprParser ts of
      ((t, []) : _) -> Right t
      _             -> Left "Parse error"
\end{minted}

O padrão \haskell{(t, [])} novamente exige que toda a
entrada tenha sido consumida.

\section{O parser de expressões com Megaparsec}

A biblioteca desenvolvida nas seções anteriores é didática, mas
bibliotecas de produção vão além: oferecem mensagens de erro
detalhadas, melhor desempenho e mais funcionalidades. O
\textbf{Megaparsec} é a biblioteca de combinadores de parsing mais
utilizada no ecossistema Haskell e segue os mesmos princípios
conceituais que estudamos, porém com uma implementação muito mais
robusta. O código discutido nesta seção está disponível em:
\begin{flushleft}
  \verb|Exp/Frontend/Lexer/Megaparsec/ExpLexer.hs| e\\
  \verb|Exp/Frontend/Parser/Megaparsec/ExpParser.hs|.
\end{flushleft}

\subsection{O tipo Parser no Megaparsec}

Nessa seção, vamos utilizar a biblioteca Megaparsec para definir um analisador
descendente recursivo para a linguagem de expressões.
Definimos o tipo \haskell{Parser} é definido como um alias para
\haskell{Parsec Void String}:

\begin{minted}{haskell}
type Parser = Parsec Void String
\end{minted}

O tipo \haskell{Parsec e s a} do Megaparsec é
parametrizado pelo tipo de \textbf{erro} \haskell|e|,
pelo tipo de \textbf{entrada} \haskell{s} e pelo tipo
do \textbf{resultado} \haskell{a}. Usando
\haskell{Void} como tipo de erro customizado,
indicamos que só utilizamos as mensagens de erro padrão da
biblioteca. \haskell{String} é o tipo da entrada:
uma lista de caracteres.

\subsection{A análise léxica com Megaparsec}

O Megaparsec separa claramente as responsabilidades léxicas das
sintáticas. O módulo

\verb|Exp.Frontend.Lexer.Megaparsec.ExpLexer|

define os primitivos léxicos que o parser irá consumir.

O primeiro componente é o \textbf{consumidor de espaço}
\haskell{sc} (\emph{space consumer}):

\begin{minted}{haskell}
sc :: Parser ()
sc = L.space space1 lineCmnt blockCmnt
 where
  lineCmnt  = L.skipLineComment "//"
  blockCmnt = L.skipBlockComment "/*" "*/"
\end{minted}

\haskell{L.space} recebe três parsers: um para
espaços em branco (\haskell{space1}, que consome um
ou mais caracteres brancos), um para comentários de linha e um para
comentários em bloco. Com isso, \haskell{sc} descarta
silenciosamente espaços e ambas as formas de comentário.

Sobre \haskell{sc}, é definido o combinador
\haskell{lexeme}, que envolve qualquer parser
consumindo automaticamente o espaço que o segue:

\begin{minted}{haskell}
lexeme :: Parser a -> Parser a
lexeme = L.lexeme sc
\end{minted}

O padrão de uso é: cada parser léxico é definido com
\haskell{lexeme}, garantindo que o espaço entre
tokens seja descartado sem que o parser sintático
precise se preocupar com isso.

Os demais primitivos léxicos são diretos:

\begin{minted}{haskell}
symbol :: String -> Parser String
symbol = L.symbol sc

parens :: Parser a -> Parser a
parens = between (symbol "(") (symbol ")")

rword :: String -> Parser ()
rword w = (lexeme . try) (string w *> notFollowedBy alphaNumChar)

int :: Parser Int
int = lexeme L.decimal
\end{minted}

\begin{itemize}
  \item \haskell{symbol} reconhece uma string literal e
    descarta o espaço seguinte.
  \item \haskell{parens p} reconhece
    \haskell{p} entre parênteses, usando
    \haskell{between} do Megaparsec para combinar os
    dois delimitadores com o conteúdo.
  \item \haskell{rword} reconhece uma \textbf{palavra
    reservada}: usa \haskell{try} para retroceder em
    caso de falha e \haskell{notFollowedBy alphaNumChar}
    para garantir que a palavra não seja prefixo de um
    identificador maior.
  \item \haskell{int} analisa um literal decimal inteiro,
    descartando o espaço seguinte.
\end{itemize}

\subsection{O analisador sintático com makeExprParser}

O analisador sintático importa do módulo
\verb|Control.Monad.Combinators.Expr| a função
\haskell{makeExprParser}, que é a versão de nível
industrial do nosso \haskell{gen} combinado com
\haskell{chainl}/\haskell{chainr}:

\begin{minted}{haskell}
import Control.Monad.Combinators.Expr
\end{minted}

A precedência e associatividade dos operadores são declaradas em uma
\textbf{tabela de operadores}, organizada do maior para o menor nível
de precedência (ou seja, os operadores que ligam mais forte vêm primeiro):

\begin{minted}{haskell}
opTable :: [[Operator Parser Exp]]
opTable
  = [ [ infixL (:*:) "*" ]
    , [ infixL (:+:) "+" ]
    ]
 where
  infixL op sym = InfixL $ op <$ symbol sym
\end{minted}

A tabela é uma lista de listas: cada lista interna agrupa operadores
de mesma precedência. \haskell{(:*:)} aparece na
primeira lista e \haskell{(:+:)} na segunda, logo
\haskell{*} tem maior precedência que
\haskell{+}. Dentro de cada grupo, os operadores têm
a mesma precedência e a associatividade é determinada pelo construtor
utilizado: \haskell{InfixL} para associatividade à esquerda.

A função auxiliar \haskell{infixL} constrói uma
entrada na tabela: \haskell{op <$ symbol sym}
reconhece o símbolo \haskell{sym} e retorna a função
construtora \haskell{op} (como \haskell{(:*:)}) como
resultado, que será aplicada pelos dois operandos adjacentes.

O parser para fatores atômicos é \haskell{termP}:

\begin{minted}{haskell}
termP :: Parser Exp
termP = parens expP
    <|> EInt <$> int
\end{minted}

\haskell{termP} tenta, em ordem: uma expressão entre
parênteses (\haskell{parens expP}) ou um literal
inteiro (\haskell{EInt <$> int}). Note que
\haskell{termP} é mutuamente recursivo com
\haskell{expP}: parênteses reiniciam o parsing a
partir do nível mais alto da gramática.

Finalmente, o parser principal é criado utilizando a tabela e o
parser \haskell{termP}:

\begin{minted}{haskell}
expP :: Parser Exp
expP = makeExprParser termP opTable
\end{minted}

\haskell{makeExprParser} recebe o parser
\haskell{termP} e a tabela de operadores, e constrói
automaticamente o parser para toda a hierarquia de precedências,
evitando a necessidade de codificar manualmente cada nível como uma
função separada.

\subsection{A função principal de parsing}

A função que executa o parser \haskell{expParser}
sobre uma \haskell{String} é:

\begin{minted}{haskell}
expParser :: String -> Either String Exp
expParser s = case parse expP "" s of
                Left err -> Left (errorBundlePretty err)
                Right e  -> Right e
\end{minted}

\haskell{parse} é a função de execução do Megaparsec:
recebe o parser, um nome de arquivo (aqui
\haskell{""}, pois não há arquivo físico) e a
entrada. Em caso de falha,
\haskell{errorBundlePretty} formata a mensagem de
erro de forma legível, com indicação de linha, coluna e o que era
esperado. Em caso de sucesso, a AST é retornada em
\haskell{Right}.

\section{Conclusão}

Este capítulo apresentou os combinadores de parsing como uma
abordagem funcional para a construção de analisadores sintáticos
recursivos descendentes. A ideia central, representar um parser
como uma função de listas de símbolos para listas de resultados,
permite instanciar o tipo \haskell{Parser} nas
classes \haskell{Functor},
\haskell{Applicative},
\haskell{Alternative} e \haskell{Monad}.

A partir dos primitivos \haskell{item} e
\haskell{sat}, construímos combinadores de alto nível
como \haskell{many},
\haskell{option}, \haskell{pack},
\haskell{chainl} e \haskell{chainr}.
Esses combinadores foram aplicados para construir um analisador
léxico e um analisador sintático para uma linguagem de expressões
aritméticas, demonstrando como a estrutura da gramática se reflete
diretamente na estrutura do código Haskell.

A combinação de \haskell{gen} com
\haskell{chainl} resolve elegantemente o problema
clássico da recursão à esquerda em gramáticas de operadores, sem a
necessidade de transformar a gramática original. O resultado é um
analisador conciso, legível e corretamente associativo que serve de
base para o frontend de um compilador.

Por fim, apresentamos o mesmo parser de expressões reescrito com o
\textbf{Megaparsec}, uma biblioteca de produção que compartilha os
mesmos fundamentos conceituais, mas oferece mensagens de erro
detalhadas (com linha, coluna e contexto), tratamento automático de
espaço em branco via \haskell{lexeme} e o combinador
\haskell{makeExprParser} para declarar hierarquias de
precedência de forma compacta através de uma tabela de operadores. A
principal diferença comportamental é a semântica de \emph{committed
choice} do operador \haskell{<|>}: ao contrário da
biblioteca desenvolvida neste capítulo, o Megaparsec só tenta a
alternativa direita se a esquerda falhou sem consumir entrada, o que
torna o parser eficiente e as mensagens de erro mais precisas.
